\chapter{Central Hypotheses of HIT}

\epigraph{\emph{A program becomes legible when its strongest claims stop hiding among its weakest ones.}}{}

The previous chapters defined the conceptual field within which HIT operates. Chapter 3 argued that harmonic information should be approached relationally rather than as a property of isolated substances, and Chapter 4 reformulated consonance as a situated function of ratio, constitution, medium, coupling regime, receptor, and context. Once those moves are in place, the framework can no longer proceed by implication alone. Its central claims must be stated explicitly, arranged by evidential status, and exposed to possible failure. This chapter therefore does not attempt to prove the whole framework in a single stroke. Its task is more exacting. It specifies what HIT currently affirms, what it treats as an open but structured hypothesis, and what remains a larger conjectural horizon.

This stratification is not rhetorical caution for its own sake. It is part of the intellectual form of the project. A framework that spans acoustics, nonlinear dynamics, neuroscience, ecology, and computational modeling cannot remain coherent if every sentence carries the same epistemic weight. Some propositions rest on relatively robust empirical footholds. Others depend on cross-domain inferences that are promising but not closed. Still others name the more ambitious consequences that would follow if the earlier layers continue to hold under stronger tests. HIT becomes more legible, not less, when those levels are distinguished rather than compressed. The aim of the present chapter is therefore to order the program without weakening it: to say clearly what kind of claim each hypothesis is, what presently supports it, and what sort of result would materially alter its standing.

\section{Why the framework requires a hierarchy of claims}

HIT requires a hierarchy of claims because the object it studies is distributed, relational, and unevenly visible across domains. Harmonic organization does not appear everywhere in the same form, at the same scale, or with the same explanatory force: what is heard as interval in one setting may appear elsewhere as a relation among modes, frequencies, couplings, or recurrent dynamical regimes. For that reason, the framework cannot be written as though one observation immediately entitled its most ambitious conclusion. The passage from structured ratio distributions to claims about cross-modal transport, informational efficiency, or biological function must be articulated in layers. This is not a sign of weakness in the framework, but part of its proper scientific form. A program that moves across acoustics, neuroscience, nonlinear dynamics, ecology, and computation becomes clearer when it distinguishes its observational footholds, its central structural hypotheses, and its larger conjectural horizon rather than compressing them into a single undifferentiated thesis.

This layered structure also protects HIT from a mistake that becomes especially likely in transdisciplinary work: treating a powerful model as though it were a literal duplicate of the world rather than one more constructed reality among others. Formal analogies across fields matter, but they do not by themselves establish that all such fields instantiate one and the same mechanism in a simple or exhaustive way \parencite{calvo_2006}. HIT therefore proceeds neither by claiming that acoustics, cortical oscillation, ecological signaling, and field dynamics are secretly identical, nor by abandoning the possibility of coordination among them. It proceeds by asking whether partially translatable patterns of ratio, recurrence, and resonant organization recur often enough, and with enough explanatory significance, to justify a common research program. Once that question is posed in a disciplined way, the architecture of the framework becomes legible: some hypotheses name empirical footholds, others define the core structural wager of the theory, and others indicate its more ambitious consequences. This makes progress easier to judge, because different parts of the program can be strengthened, limited, or revised without forcing the whole framework into either premature closure or shapeless generality.

The six hypotheses can now be stated in their basic order before being developed one by one: \textbf{H1}, natural signals contain structured ratio distributions; \textbf{H2}, these distributions are learnable by computational systems; \textbf{H3}, ratio structure is shared cross-modally; \textbf{H4}, simple harmonic ratios correspond to lower processing cost or greater informational efficiency; \textbf{H5}, biological systems possess a functional sensitivity to consonant organization; and \textbf{H6}, the interval or ratio can function as a scale-invariant carrier of information.

\section{H1 and H2: the empirical footholds}

The first two hypotheses provide the framework's most stable empirical footing. \textbf{H1} states that natural signals contain structured ratio distributions that are not well described as random residue. \textbf{H2} states that such distributions are learnable by computational systems. These are modest claims relative to the larger ambitions of HIT, but they are indispensable. Without them, harmonic information would remain a suggestive vocabulary in search of an object.

\textbf{H1} belongs to the observational level. The point is not that every signal displays obvious low-order harmonic regularity at every scale. The point is that, across many natural regimes, ratio structure departs from arbitrariness strongly enough to support lawful description. Recurrence-based work already points in that direction. Trulla and colleagues showed that simpler harmonic relations leave measurable signatures in temporal organization, suggesting that ratio structure has dynamical consequence rather than retrospective descriptive convenience \parencite{trulla_2018}. The wider literatures surveyed earlier push in the same direction from other angles: neural entrainment, cross-frequency coupling, harmonicity models, ecological partitioning, and mode organization all imply that structured proportional relations recur where systems sustain coordinated behavior.

\textbf{H2} adds a computational test. Within Phideus (the project's computational probe, developed to test harmonic descriptors and cross-domain learning, and discussed more fully in Chapter 11), ratio-structured descriptors and ratio-sensitive representations have repeatedly proved learnable in ways that affect signal organization and, under some conditions, downstream alignment. That is stronger than the generic observation that flexible models can fit data. Learnability does not settle ontology, but it is a strong filter against empty projection. A structure that cannot be learned at all is unlikely to support the kind of explanatory program HIT requires.

Taken together, \textbf{H1} and \textbf{H2} justify a dedicated research program while leaving the harder questions open. They do not yet specify which ratio organizations matter most, how far the structuredness extends beyond the cases studied so far, or whether the learned structure is causally privileged. They establish something more basic and more important at this stage: there is a tractable object here, and it is patterned enough to deserve systematic investigation.

\section{H3 and H6: shared structure and invariant carriers}

The next layer contains the hypotheses that give HIT its distinct conceptual profile. \textbf{H3} states that ratio structure is shared cross-modally. \textbf{H6} states that the interval, or ratio more generally, can function as a scale-invariant carrier of information. The two are linked but not interchangeable. \textbf{H3} concerns the recurrence of relational structure across domains. \textbf{H6} concerns the status of the ratio itself as a candidate informational primitive.

The formal core of \textbf{H6} can be stated very simply. Given two frequencies or modes, the relevant interval may be written as $r = f_2 / f_1$. When the relation is reduced to lowest integers, one may write $r_{\mathrm{red}} = p / q$. If both terms are rescaled by a common constant $k$, so that $f'_1 = k f_1$ and $f'_2 = k f_2$, then $r' = f'_2 / f'_1 = r$. This does not imply that every medium or receptor will respond identically across scales. It isolates a narrower point: the proportional relation survives absolute rescaling, which is precisely why the ratio is a plausible candidate carrier of structure rather than a local accident of pitch or magnitude.

\textbf{H3} remains under active test, but it is no longer a purely speculative proposition. The strongest support currently comes from the computational side of the program, where cross-modal descriptor-guided learning has shown that ratio-sensitive structure can reorganize latent geometry and improve retrieval under controlled conditions. That is not a universal theorem of modality-sharing. It is enough, however, to show that harmonic organization need not remain confined to a single representational substrate. Neuroscientific work on cross-frequency coupling points in a compatible direction, since relations among oscillatory scales already operate there as integrative rather than merely local features \parencite{canolty_2010}. The experimental burden now is comparative: the transfer must survive matched controls and interventions, not just descriptive fit.

\textbf{H6} reaches deeper because it asks what kind of thing could travel in the first place. In Shannon's framework, information is defined over possible messages rather than over one privileged material realization \parencite{shannon_1948}. HIT extends that intuition into a relational register: under some conditions, a ratio may carry organization because what matters is the proportional relation itself. Kuramoto-class dynamics make the idea concrete. Simple ratios repeatedly emerge as privileged loci of locking, coordination, and stability across systems with very different material realizations \parencite{kuramoto_1984}. Work on traveling waves and cortical geometry suggests a comparable possibility: eigenfrequency relations may encode organizational properties that are not reducible to one local measurement alone \parencite{jacobs_2025}.

At this point the trans-physical framing becomes decisive. The interval is not being proposed as a mystical token hidden inside all phenomena. It is being proposed as a particularly promising relational form. In acoustics, that form appears as an audible interval. In neural dynamics, it may appear as coupling among oscillatory bands. In other physical settings, it may appear as a relation among modes, resonant frequencies, or recurring scales of organization. What matters is not literal sameness across those cases, but whether the ratio preserves enough structural identity to function as a common analytic object. That is the wager shared by \textbf{H3} and \textbf{H6}: structured proportional relations may be more portable, and more explanatory, than standard domain partitions have allowed.

\section{H4 and H5: the major conjectures}

The last layer contains the framework's most ambitious claims. \textbf{H4} states that simple harmonic ratios minimize processing cost, or at least correspond to regimes of greater informational efficiency. \textbf{H5} states that biological systems possess a functional sense of consonance. These propositions reach beyond structure, learnability, and cross-modal organization toward economy, viability, and adaptive significance.

\textbf{H4} is best understood as a structured conjecture. Landauer's principle established that information is not free from physical cost \parencite{landauer_1961}. Later work on the thermodynamics of prediction argued that structured regularities can reduce the burden of inference when they are effectively exploited \parencite{still_2012}. Friston's free-energy framework extends a related intuition into adaptive systems, where viability depends in part on attunement to lawful environmental structure \parencite{friston_2010}. Taken together, these lines make it reasonable to ask whether regular, recurrent, low-conflict relations are often cheaper to predict, encode, or stabilize than irregular alternatives. What would strengthen the hypothesis is an operational link among ratio simplicity, recurrence profile, prediction cost, and energetic or computational expenditure. That link is still being built.

\textbf{H5} is ambitious in a different register. It proposes that living systems may possess a functional sensitivity to consonant organization, where consonance names stable, interpretable, low-conflict coordination rather than a merely aesthetic preference. Neural resonance research suggests that the nervous system is not indifferent to patterned harmonic relation \parencite{large_2012}. Reward-related studies indicate that pleasurable musical experience can couple to affective and reward-related neural processes under some conditions \parencite{blood_2001}. Krause's acoustic ecology and Glass and Mackey's work on physiological ratio-locking widen the frame beyond human listening, while Levin's work on bioelectric organization shows that living form can depend on patterned ratios and gradients rather than only on local molecular identity \parencite{glass_1988, krause_2013, levin_2021}.

Recent work sharpens this conjecture from a more literal physiological side. Kawai reports slow harmonic cross-frequency coupling among respiratory, cardiac, and brainstem oscillators in vivo, while Zheng and colleagues show that near-harmonic alpha-theta relations track attentional organization in human EEG \parencite{kawai_2023, zheng_2025}. Bahuguna and colleagues extend the point into high-dimensional oscillatory portraits whose interdependence patterns predict visuomotor behavior, and Medvedev and Lehmann show that cross-frequency coupling matrices can support high-performance classification of absence seizures under deep learning pipelines \parencite{bahuguna_2025, medvedev_2025}. In all four cases, what becomes informative is not one band or one signal alone, but the patterned relation among bands or rhythms.

The critical point is one of scope. A functional sense of consonance, if it exists, need not appear as conscious musical preference. It may emerge instead as sensitivity to coherent relational organization, to regimes of coupling that preserve viability, interpretability, or coordinated action. This makes the conjecture both broader and more testable. The question is not whether every organism \enquote{likes} consonance. It is whether biological systems detect, exploit, or depend on low-conflict harmonic organization in ways that materially affect function.

Together, \textbf{H4} and \textbf{H5} define the horizon toward which HIT is moving. They are the reason the framework does not stop at description, learning, and transfer. For now, they are best held as informed conjectures with a growing conceptual and empirical perimeter rather than as established laws. Their value lies partly in the fact that they can be sharpened, limited, or rejected by the right experiments.

\section{Status, falsifiers, and the next decisive tests}

Table~\ref{tab:hypotheses} gathers the six hypotheses into a compact view that situates each claim within the program's evidential architecture.

\begin{table}[ht]
\centering
\caption{Central hypotheses of HIT, ordered by evidential status.}
\label{tab:hypotheses}
\footnotesize
\begin{tabular}{@{}cp{4.2cm}lp{3.0cm}@{}}
\toprule
 & Hypothesis & Status & Strongest evidence \\
\midrule
H1 & Structured ratio distributions exist & Empirical foothold & Learnable by NNs \\
H2 & Ratio distributions are learnable & Empirical foothold & Phideus: $\text{val\_loss} < 0.5$ \\
H3 & Ratio structure is shared cross-modally & Central wager & Escalon~1: $S = 84.1\%$ \\
H4 & Simple ratios minimize processing cost & Conjecture & Recurrence analysis \\
H5 & Biological systems sense consonance & Conjecture & Neural resonance lit. \\
H6 & Ratio is a scale-invariant carrier & Central wager & Kuramoto dynamics \\
\bottomrule
\end{tabular}
\end{table}

The architecture of the chapter can now be gathered into a more compact view. \textbf{H1} and \textbf{H2} are the program's footholds: they state that structured ratio distributions exist in natural signals and can be computationally learned. \textbf{H3} and \textbf{H6} define its central structural wager: ratio relations may travel across modalities and operate as scale-robust carriers of organization. \textbf{H4} and \textbf{H5} mark the framework's larger consequences: harmonic simplicity may be tied to informational efficiency, and living systems may exhibit functional sensitivity to consonant regimes. These hypotheses differ in strength, but they also depend on one another asymmetrically. If the lower layers fail, the upper ones lose support. If the upper ones contract, the lower ones may still remain intact.

Each hypothesis also has a distinct exposure to falsation or revision. \textbf{H1} would be weakened if structured ratio distributions consistently dissolved under better controls or broader datasets. \textbf{H2} would lose force if ratio-sensitive learning repeatedly failed in settings where the structure is supposed to matter. \textbf{H3} depends on strict comparative protocols: if cross-modal sharing disappears under causal controls, its strongest formulations cannot survive. \textbf{H6} would be pressured if interval-like relations turned out to be far less portable than the framework assumes, or if scale invariance repeatedly collapsed outside narrow acoustic cases. \textbf{H4} stands or falls with operational links among structure, predictability, and cost. \textbf{H5} depends on whether biological systems can be shown to register harmonic organization functionally rather than only aesthetically or retrospectively.

This means that the experimental and theoretical future of HIT is not vague. It is already structured by the hypotheses themselves. Some tests belong to representation learning and causal comparison. Others belong to dynamical modeling, recurrence analysis, psychophysics, neurophysiology, or embodied interaction. Still others belong to theory construction, where the relation among information, resonance, and efficiency must be stated more sharply than it currently is. What makes the framework viable is precisely that these tests are not interchangeable. Different outcomes would reshape different parts of the map. A null result for one conjecture would not automatically invalidate the whole program; a strong positive result in one domain would not automatically settle every other.

The advantage of this architecture is therefore twofold. It preserves ambition where ambition is warranted, and it preserves discipline where discipline is indispensable. HIT is not a single declaration that the world is harmonic in some undifferentiated sense. It is a structured proposal about patterned relation, learnability, shared organization, efficiency, and function. Some of these elements now stand on firmer ground than others. That is as it should be. A serious research program is defined not by flattening those differences, but by making them explicit.

The next task follows from this directly. Once the hypotheses have been stated and ranked, they must be set against the actual state of knowledge. We need to know which domains have already uncovered relevant parts of the picture, which vocabularies they use, and what remains invisible from within their local frames. That broader map belongs to the next part of the manuscript.
